\chapter{State of the Art}

\section*{Introduction}

This chapter examines existing monitoring solutions, available anomaly detection
techniques, and recent advances in language models applied to IT operations.
This comparative analysis is essential to justify the architectural choices of
\caimp{} and demonstrate the added value of the adopted approach.

\section{Commercial Monitoring Solutions}

\subsection{Datadog}

Datadog is a SaaS monitoring platform founded in 2010, with over 27,000 clients
in 2023. Its \textit{Watchdog} AI feature provides automatic anomaly detection,
but explanations remain superficial and do not leverage a generative LLM.
Pricing: \$15--31/host/month; no on-premise deployment available.

\subsection{Grafana + Prometheus}

The Grafana/Prometheus combination has become the open source observability
standard. Prometheus uses a pull-based model with its own TSDB and PromQL query
language. Grafana provides visualisation across dozens of data sources. Both are
entirely open source, but have \textbf{no native AI} and require significant
expertise to configure. Alert fatigue is common without intelligent deduplication.

\subsection{Dynatrace}

Dynatrace features the Davis AI engine, a causal analysis system that
automatically correlates anomalies across the full application stack. The most
advanced root cause analysis on the market, but at \$69/host/month and with no
open-source or on-premise option.

\subsection{Splunk + ITSI}

Splunk Enterprise with ITSI \cite{splunk_itsi} adds MLTK-based anomaly detection
and event correlation. \caimp{} provides optional integration with Splunk as a
secondary metric source. Already-equipped organisations can leverage \caimp{}'s
explainability layer on top of Splunk's detection.

\subsection{Open Source Alternatives}

SigNoz (built on ClickHouse), Netdata (real-time single-host), Zabbix (mature
but steep learning curve), and Nagios (static threshold pioneer, now outdated)
all lack native AI explanation capabilities.

\section{Comparative Table}

\begin{table}[H]
\centering
\caption{Comparison of monitoring solutions across 12 criteria}
\label{tab:comparison}
\small
\begin{tabularx}{\textwidth}{lXXXXXX}
\toprule
\textbf{Criterion} & \textbf{Datadog} & \textbf{Grafana/Prom.} & \textbf{Dynatrace}
  & \textbf{SigNoz} & \textbf{Zabbix} & \textbf{\caimp{}} \\
\midrule
Price/host/mo  & \$15--31    & Free       & \$69     & Free        & Free       & Free \\
On-premise     & No          & Yes        & No       & Yes         & Yes        & Yes \\
Native AI      & Partial     & No         & Advanced & No          & No         & Yes \\
LLM explanation& No          & No         & No       & No          & No         & \textbf{Yes} \\
RAG + runbooks & No          & No         & No       & No          & No         & \textbf{Yes} \\
Change corr.   & No          & No         & Partial  & No          & No         & \textbf{Yes} \\
24h forecast   & Yes         & No         & Yes      & No          & No         & Yes \\
Alert dedup.   & Yes         & Via AM     & Yes      & Partial     & No         & Yes \\
Multi-tenant   & Yes         & Limited    & Yes      & Yes         & Yes        & Yes \\
Deploy $<$10m  & No          & No         & No       & No          & No         & \textbf{Yes} \\
Answers-first  & No          & No         & Partial  & No          & No         & \textbf{Yes} \\
\bottomrule
\end{tabularx}
\end{table}

\section{Anomaly Detection Techniques}

\subsection{Static Threshold}

The simplest method: alert when a value exceeds a predefined threshold.
Fast and deterministic, but prone to false positives on variable-load systems.
Implemented in \caimp{} as a first rapid filter.

\subsection{Z-Score Statistical Detection}

The Z-score measures how many standard deviations a reading is from the
historical mean (Equation~\ref{eq:zscore}):
\begin{equation}
  z = \frac{x - \mu}{\sigma}
  \label{eq:zscore}
\end{equation}
A value $|z| > 3$ is considered anomalous. This method adapts automatically
to the system's normal behaviour.

\subsection{Rate of Change}

The discrete derivative of the time series:
\begin{equation}
  \Delta_t = \frac{x_t - x_{t-1}}{x_{t-1}} \times 100\%
\end{equation}
Particularly effective for detecting rapid saturations that Z-score on a long
window might miss.

\subsection{Holt-Winters Double Exponential Smoothing}

\cite{holt1957forecasting,winters1960forecasting}:
\begin{align}
  L_t &= \alpha x_t + (1-\alpha)(L_{t-1} + T_{t-1}) \\
  T_t &= \beta(L_t - L_{t-1}) + (1-\beta)T_{t-1} \\
  \hat{x}_{t+h} &= L_t + h \cdot T_t
\end{align}
Parameters used in \caimp{}: $\alpha = 0.3$, $\beta = 0.1$.
Confidence intervals: $\pm 1.28\sigma$ residuals ($\approx$80\% coverage).

\section{Language Models for AIOps}

\subsection{Open LLM Landscape}

\textbf{Llama 3.1 8B} \cite{llama3_meta2024} was selected for \caimp{} for its
excellent quality/performance ratio in Q4\_K\_M quantised format: 4.7\,GB,
inference latency under 2 seconds on standard hardware. Alternatives evaluated:
Mistral 7B (slightly larger context window), Phi-3 Mini (excellent on constrained
hardware), Gemma 7B (less well-aligned for operational tasks).

\subsection{Quantisation}

Q4\_K\_M quantisation \cite{dettmers2022quantization} reduces Llama 3.1 8B from
16\,GB (fp16) to 4.7\,GB with less than 2\% quality degradation on reasoning
benchmarks, making local LLM inference feasible on standard server hardware.

\subsection{RAG Pattern}

The Retrieval-Augmented Generation pattern \cite{lewis2020rag}:
\begin{enumerate}
  \item \textbf{Indexing}: documents (runbooks, past incidents) are embedded via
    \texttt{nomic-embed-text} and stored in pgvector \cite{pgvector2021}.
  \item \textbf{Retrieval}: at incident time, a query embedding is computed and
    the $k$ most similar documents are retrieved via HNSW cosine similarity search.
  \item \textbf{Augmentation}: retrieved documents are injected into the LLM
    prompt as context before generation.
\end{enumerate}

\section{Positioning of CAIMP}

The state of the art reveals a gap between expensive commercial AIOps solutions
(Datadog, Dynatrace) and open source tools without AI (Grafana/Prometheus).
\caimp{} occupies this space with four key differentiators: natural-language
explanation via a local LLM; automatic change correlation (absent from all
compared open source solutions); a lightweight modern event-driven architecture;
and ten-minute deployment on any machine.

\section*{Conclusion}

The state of the art shows that while monitoring tools are plentiful, none
simultaneously combines self-hosting, accessibility for non-experts, and natural
language explanation via a local LLM. This gap fully justifies the creation of
\caimp{}. The following chapter translates these findings into concrete functional
and non-functional specifications.
